\section{Where the method failed}

Four of the seven mechanisms in \S4 arrived late, after the failure they were built for had already happened. That is not damning --- a method extracted from one project will be, by construction, a record of what that project learned the hard way. Those are covered where they occurred.

This section is about the failures that are still live. The last one is the serious one, and I want to arrive at it properly.

\subsection{Three that arrived late}

\textbf{The inherited marker was applied to numbers and not to methods.} Day one for figures, experiment twelve for criteria --- by which point a carried statistical bar had governed ten experiments. The asymmetry is the finding: a repository that would not let a number cross without provenance let a criterion cross unmarked, and the criterion determined how every number was read.

\textbf{Verification ran per passage while repetition crossed passages.} Found only because the same ratio appeared twice and someone happened to look at both.

\textbf{Specification defects had no ledger.} Decisions got foreclosures, measurements got provenance, and badly-written specifications got nothing --- so the same two errors recurred: bundling a bounded task with an unbounded one, and naming a branch without asking for a commit. Both were cheap and both were avoidable by a five-line list that did not exist.

All three now ship in the template, each with the failure that motivated it. That is the ordinary way a method matures, and it is unremarkable.

\subsection{The devil's advocate was agreed and never written down}

Less ordinary, and worth its own paragraph because of what it demonstrates.

The pass was proposed, discussed, refined into two triggers, and endorsed. It ran twice, informally, and both times produced something the loop would not have produced. \textbf{It was never written into the workflow file.}

By the method's own foundational rule --- every handoff is a versioned artifact, never chat context --- that means it did not exist. A mechanism designed to catch the failure of a method going unquestioned was itself never subjected to the method. If the sessions had been cleared, it would have vanished without trace, and nobody would have known to miss it.

\subsection{And now the one that matters}

Two components of the method were never used. Not underused. \textbf{Zero times.}

The low-ceremony lane --- where an idea can be tried and thrown away, no decision record, no findings entry, no definition of done, one permitted output --- shipped on the first day with its rules, its README, and a test that enforces its one constraint in continuous integration. It was used zero times in thirteen experiments.

The third surface, for cross-experiment synthesis and work spanning many artifacts, was designed into the workflow, given a role in the table, and used zero times. The synthesis it existed for happened in a chat window and had to be reconstructed afterwards from recollection.

The obvious reading is oversight. It is not, and the evidence against it is specific and damning.

\textbf{The lane was built as a fix for exactly this failure, in the previous project, where it had already cost something.} In that project a step was written up twice and never started, accumulating five declared items, four of them open questions, because the constitution had one standard --- plan, decision record, findings, tests, definition of done --- and that standard is right for library code and far too expensive for exploration. When a question is open and the only sanctioned move is to write about it, uncertainty turns into documents. The diagnosis was made explicitly. The remedy was designed explicitly. It was carried into the new repository on day one, complete with enforcement.

And then the new project \textbf{did the same thing again.} Not by forgetting the lane, but by never reaching for it, while writing preambles and specifications and analyses about questions that could have been settled in twenty minutes by trying them.

So the finding is not \textit{the method lacked a lane for informal work.} It is:

\begin{quote}
\textbf{A project with strong ceremony does not reach for its low-ceremony lane, even when the lane exists, is enforced, is documented, and was built specifically because its absence had already caused harm.}
\end{quote}

The mechanism is not mysterious once stated. Every artifact that meets a high standard raises the felt cost of producing one that does not. After ten pre-registered experiments with falsifiers, verdicts and provenance, writing a directory containing a rough script and no record does not feel like the sanctioned move it formally is. \textbf{It feels like a lapse.} The rigour is self-reinforcing, and self-reinforcement does not distinguish between the parts of itself that are load-bearing and the parts that are ritual.

This has a second face which is worse. The formal parts of the method are the ones that produce artifacts, and artifacts are the only thing the method can see. A spike that would have settled a question in an afternoon leaves nothing behind if it succeeds --- the whole point is that it is deleted. \textbf{So the method is structurally blind to the value of the work it is discouraging}, and the record it produces will always show that the ceremony was worth it, because the record is made of ceremony.

I felt this directly and did not name it at the time. There were several afternoons of writing about a question --- carefully, with pointers and open questions and a plan for resolving it --- where the honest move was to spend twenty minutes finding out. The writing was good. It met the standard. It also postponed the answer, and it was easier than admitting I did not know, because a document that says \textit{here are the four open questions} looks like progress and a directory called \texttt{does-the-fixture-need-occlusion} looks like not knowing.

That is the thing a methodology cannot fix by adding a rule, which is why it is the last section rather than a bullet in \S4. \textbf{Adding the directory was not enough. It was already there.}

\subsection{What might actually work}

Not confident about this, and stating it as a proposal rather than a finding.

A \textbf{trigger} rather than a lane, in the same shape as the devil's advocate pass: \textit{when a question has been written about twice without being tested, it goes to a spike.} Twice is countable. It fires on the observable --- documents produced --- rather than on a judgement about whether something needs trying, which is exactly the judgement the ceremony corrupts.

And a \textbf{default toward the cheap move at the moment of maximum uncertainty}. The method's autonomy ladder governs how much the collaborator decides; nothing governs how much rigour the artifact must carry, and the answer should vary. At the point where a question is genuinely open, the standard should drop, not rise.

Both are in the template. Neither has been tested, and the reader should treat them accordingly --- as the most recent guesses of someone who has now made this mistake twice.
